% ----------------------------------------------------------
% ELEMENTOS PRÉ-TEXTUAIS
% ----------------------------------------------------------
% \pretextual

% ---
% Capa
% ---
\imprimircapa
% ---

% ---
% Folha de rosto
% (o * indica que haverá a ficha bibliográfica)
% ---
%\imprimirfolhaderosto*
\imprimirfolhaderosto
% ---



% ---
% Agradecimentos
% ---
\begin{agradecimentos}
Agradeço aos meus pais, Simone e Joel, pelo amor incondicional, pela educação que me
formou como pessoa e pelo apoio em cada etapa da minha trajetória acadêmica. Sem a base
que construíram, este trabalho não seria possível.

À minha namorada, Karolline, pelo companheirismo constante, pela paciência nos momentos
de ausência e pelo incentivo que nunca faltou nos momentos mais difíceis deste trabalho.
Sua presença tornou essa jornada mais leve e significativa.

Ao meu orientador, Jean Schmith, pela disponibilidade, pelas orientações precisas e pela confiança
depositada neste trabalho.

Aos professores do curso de Engenharia da Computação da UNISINOS, pelos ensinamentos que
fundamentam esta e todas as demais etapas da minha formação.

A todos que, de alguma forma, contribuíram para que eu chegasse até aqui, meu sincero
agradecimento.
\end{agradecimentos}
% ---


% ---
% RESUMOS
% ---

% resumo em português
\setlength{\absparsep}{18pt} % ajusta o espaçamento dos parágrafos do resumo
\begin{resumo}
Investigou-se a caracterização de padrões de escoamento bifásico com isobutano em
microtubo horizontal de 1\,mm por meio de sinais de intensidade extraídos de panoramas
de imagens. Analisaram-se 26 ensaios de cinco padrões, formando 87 pares do mesmo padrão
e 238 de padrões diferentes. Extraiu-se a intensidade média por coluna em uma faixa dos panoramas.
Os sinais foram invertidos, tiveram a tendência linear removida e foram normalizados. Aplicou-se a transformada
\textit{wavelet} contínua de Morlet às funções de correlação para obter os
espectros de potência e cruzado e a coerência
wavelet alternativa. Comparou-se diretamente a forma dos perfis de potência no atraso
zero dos panoramas completos pelo coeficiente de Pearson $r$. Os mapas de coerência wavelet
utilizaram os trechos iniciais de comprimento comum dos sinais.
Examinaram-se visualmente exemplos selecionados e resumiu-se cada mapa por seus
valores mínimos, médios, medianos e máximos.
Compararam-se também as coerências wavelet médias dos pares dentro e entre padrões.
O coeficiente de Pearson $r$
apresentou média de $0{,}7642$ dentro dos padrões e de $0{,}5814$ entre padrões,
indicando maior semelhança espectral média interna, acompanhada de sobreposição das
distribuições. A média das coerências médias dos pares foi de $0{,}0607$
dentro dos padrões e de $0{,}0874$ entre padrões, apresentando comportamento
distinto de Pearson. Em análise complementar, estimaram-se a
fração de vazio e a densidade espacial de transições a partir de máscaras de gás,
nas mesmas 54 linhas utilizadas na análise wavelet. Foram elegíveis 25 ensaios, após a exclusão
de um registro saturado. O anular apresentou a maior mediana da fração de vazio
estimada e a menor mediana da densidade de transições. Essas medidas apoiaram a interpretação
dos descritores espectrais. A estimativa de fração de vazio permaneceu sem calibração
volumétrica independente. Os pares compartilharam aquisições e os grupos tiveram números
diferentes de ensaios. A avaliação permaneceu descritiva, sem testes de
significância estatística ou estimativa de acurácia de classificação.

 \textbf{Palavras-chave}: escoamento bifásico. microcanais. transformada \textit{wavelet} contínua. coerência \textit{wavelet}. fração de vazio. identificação de padrões de escoamento.
\end{resumo}

% resumo em inglês
\begin{resumo}[Abstract]
 \begin{otherlanguage*}{english}
Two-phase flow patterns of isobutane in a horizontal 1\,mm microtube were investigated
using intensity signals extracted from image panoramas. Twenty-six acquisitions from
five patterns were analyzed, yielding 87 within-pattern and 238 between-pattern pairs.
The mean intensity per column within a panorama band was extracted, followed by intensity inversion,
linear detrending and normalization. The Morlet continuous wavelet transform was applied to correlation
functions to obtain the power and cross
spectra and alternative wavelet coherence. The shapes of zero-lag power profiles from full
panoramas were compared directly using Pearson's correlation coefficient $r$. Wavelet
coherence maps used equal-length initial segments of the signals.
Selected examples were examined visually, and every map was summarized by its
minimum, mean, median, and maximum values.
Mean wavelet coherence per pair was also compared within and between patterns.
Mean Pearson's correlation coefficient $r$ was 0.7642 within patterns and 0.5814 between
patterns, indicating greater mean within-pattern spectral similarity, with overlapping
distributions. The average of pairwise mean coherence values was 0.0607 within
patterns and 0.0874 between patterns, showing a different pattern from Pearson's
correlation. A complementary analysis estimated void fraction and
spatial transition density from gas masks over the same 54 rows used in the wavelet analysis.
Twenty-five acquisitions were eligible after excluding one saturated record.
The annular pattern exhibited the highest median estimated void fraction and the lowest
median transition density. These measures supported the interpretation of
spectral descriptors. The void fraction estimate lacked independent volumetric calibration.
Pairs shared acquisitions, and groups had different numbers of acquisitions.
The evaluation remained descriptive, without statistical significance testing
or estimation of classification accuracy.

   \vspace{\onelineskip}

   \noindent
   \textbf{Keywords}: two-phase flow. microchannels. continuous wavelet transform. wavelet coherence. void fraction. flow pattern identification.
 \end{otherlanguage*}
\end{resumo}


% ---
% inserir lista de ilustrações
% ---
\pdfbookmark[0]{\listfigurename}{lof}
\listoffigures*
\cleardoublepage
% ---

% ---
% inserir lista de quadros
% ---
\pdfbookmark[0]{\listofquadrosname}{loq}
\listofquadros*
\cleardoublepage
% ---

% ---
% inserir lista de tabelas
% ---
\pdfbookmark[0]{\listtablename}{lot}
\listoftables*
\cleardoublepage
% ---

% ---
% inserir lista de abreviaturas e siglas
% ---
\begin{siglas}
  \item[CNN] Rede Neural Convolucional (\textit{Convolutional Neural Network})
  \item[CWT] Transformada Wavelet Contínua (\textit{Continuous Wavelet Transform})
  \item[DWT] Transformada Wavelet Discreta (\textit{Discrete Wavelet Transform})
  \item[FFT] Transformada Rápida de Fourier (\textit{Fast Fourier Transform})
  \item[LDR] Resistor Dependente de Luz (\textit{Light Dependent Resistor})
  \item[LED] Diodo Emissor de Luz (\textit{Light Emitting Diode})
  \item[LETEF] Laboratório de Estudos Térmicos e Fluidodinâmicos
  \item[LWE] Coeficiente de Energia Wavelet Local (\textit{Local Wavelet Energy coefficient})
  \item[OS] Sensor Óptico (\textit{Optical Sensor})
  \item[PD] Transdutor de Pressão Diferencial
  \item[PH] Pré-aquecedor (\textit{Pre-Heater})
  \item[PSD] Densidade Espectral de Potência (\textit{Power Spectral Density})
  \item[SNR] Relação Sinal-Ruído (\textit{Signal-to-Noise Ratio})
  \item[TS] Seção de Teste (\textit{Test Section})
  \item[VS] Seção de Visualização (\textit{Visualization Section})
  \item[WCO] Coerência Wavelet (\textit{Wavelet Coherence})
  \item[WCS] Espectro Cruzado Wavelet (\textit{Wavelet Cross Spectrum})
  \item[WPS] Espectro de Potência Wavelet (\textit{Wavelet Power Spectrum})
\end{siglas}
% ---

% ---
% inserir lista de símbolos
% ---
\begin{simbolos}
  % --- Quantidades espectrais ---
  \item[$W_x(a,b)$] Coeficiente da CWT do sinal $x$ na escala $a$ e posição $b$
  \item[$P_x(a,b)$] Espectro de potência wavelet clássico do sinal $x$
  \item[$C_{xy}(a,b)$] Espectro cruzado wavelet clássico entre os sinais $x$ e $y$
  \item[$\gamma_{xy}(a,b)$] Coerência wavelet clássica quadrática entre os sinais $x$ e $y$
  \item[$R_{xx}(\tau)$] Função de autocorrelação do sinal $x$
  \item[$R_{xy}(\tau)$] Função de correlação cruzada entre os sinais $x$ e $y$
  \item[$\widetilde{P}_x(a,b)$] Campo de potência wavelet alternativo do sinal $x$
  \item[$\widetilde{C}_{xy}(a,b)$] Espectro cruzado wavelet alternativo entre os sinais $x$ e $y$
  \item[$\widetilde{\gamma}_{xy}(a,b)$] Coerência wavelet alternativa entre os sinais $x$ e $y$
  \item[$\lambda_k$] Comprimento de onda espacial na escala $k$ [px]
  \item[$P_{ik}$] Potência alternativa no atraso nulo do ensaio $i$ na escala $k$
  \item[$M_{ik}$] Potência alternativa média nos atrasos do ensaio $i$ na escala $k$
  \item[$r_{ij}$] Coeficiente de Pearson entre os perfis de potência dos ensaios $i$ e $j$
  % --- Descritores ópticos complementares ---
  \item[$\alpha$] Fração volumétrica de vazio
  \item[$\widehat\alpha_{\mathrm{opt}}$] Estimativa óptica condicional de fração de vazio, média nas colunas
  \item[$\widehat\alpha_{\mathrm{opt},j}$] Fração de vazio estimada na coluna $j$
  \item[$d_{ij}$] Largura radial do intervalo de gás $i$ na coluna $j$ [px]
  \item[$D_{\mathrm{px}}$] Diâmetro interno adotado na análise das máscaras [px]
  \item[$H$, $W$] Número de linhas e número de colunas da máscara, respectivamente
  \item[$N_{\mathrm{trans}}$] Número de mudanças de valor entre pixels vizinhos na mesma linha da máscara
  \item[$F$] Densidade espacial de transições da máscara por par de pixels vizinhos
  % --- Variáveis termofluidodinâmicas ---
  \item[$x$] Título de vapor, razão entre a vazão mássica de vapor e a vazão mássica total
  \item[$S$] Razão de deslizamento entre as velocidades médias do vapor e do líquido
  \item[$\beta$] Parâmetro adimensional de identificação do padrão de escoamento
  \item[$G$] Fluxo mássico [kg\,m$^{-2}$\,s$^{-1}$]
  \item[$q''$] Fluxo de calor [kW\,m$^{-2}$]
  \item[$Co$] Número de confinamento
  \item[$\sigma$] Tensão superficial [N\,m$^{-1}$]
  \item[$\rho_l$] Massa específica do líquido [kg\,m$^{-3}$]
  \item[$\rho_v$] Massa específica do vapor [kg\,m$^{-3}$]
  \item[$D_h$] Diâmetro hidráulico [m]
  \item[$i_{lv}$] Calor latente de vaporização [J\,kg$^{-1}$]
\end{simbolos}
% ---

% ---
% inserir o sumario
% ---
\pdfbookmark[0]{\contentsname}{toc}
\tableofcontents*
\cleardoublepage
% ---
